\documentclass[11pt,a4paper]{article}

\usepackage[utf8]{inputenc}
\usepackage[T1]{fontenc}
\usepackage{lmodern}
\usepackage[margin=2.5cm]{geometry}
\usepackage{amsmath,amssymb,amsfonts}
\usepackage{booktabs}
\usepackage{xcolor}
\usepackage{enumitem}
\usepackage{titlesec}
\usepackage{fancyhdr}
\usepackage[most]{tcolorbox}
\usepackage[round,authoryear]{natbib}
\usepackage[hidelinks,colorlinks=true,linkcolor=blue!70!black,citecolor=blue!70!black,urlcolor=blue!70!black]{hyperref}

% Page styling
\setlength{\headheight}{14pt}
\pagestyle{fancy}
\fancyhf{}
\lhead{\textsc{Response to Reviewer 2 --- Dr.~Antoine Deleforge}}
\rhead{\textsc{Gerardo Cicalese}}
\cfoot{\thepage}

\definecolor{commentbg}{RGB}{245, 247, 250}
\definecolor{commentframe}{RGB}{70, 95, 130}
\definecolor{responseframe}{RGB}{40, 110, 70}

\newtcolorbox{revcomment}[2][]{
    enhanced,
    breakable,
    colback=commentbg,
    colframe=commentframe,
    fonttitle=\bfseries,
    title={#2},
    boxrule=0.8pt,
    arc=2mm,
    top=2mm, bottom=2mm, left=3mm, right=3mm,
    before skip=3mm, after skip=2mm,
    #1
}

\newenvironment{response}{
    \par\vspace{1mm}\noindent\textbf{\textcolor{responseframe}{Response:}}\par\vspace{1mm}
}{
    \par\vspace{3mm}
}

\begin{document}

\title{\textbf{\Large Formal Response to Reviewer 2 (Dr.~Antoine Deleforge)}\\[0.3em]
\large Re: Ph.D. Dissertation Review\\\emph{``Domain Decomposition and Differentiable Physics for the Viscoelastic Wave Equation''}}
\author{\textbf{Candidate:} Gerardo Cicalese\\[0.3em]
\textbf{Advisor:} Prof.~Ilario Mazzieri\\[0.3em]
\textbf{Co-Advisors:} Prof.~Gabriele Ciaramella, Prof.~Fabio Antonacci\\[0.3em]
\textbf{Institution:} Politecnico di Milano}
\date{November 2026}
\maketitle

\noindent Dear Dr.~Antoine Deleforge,\\[0.5em]
We express our deepest gratitude for your comprehensive, constructive, and rigorous review of the doctoral dissertation \emph{``Domain Decomposition and Differentiable Physics for the Viscoelastic Wave Equation''}. Your critical feedback has pinpointed key areas requiring physical nuance, formal mathematical correction, toned-down scientific claims, and substantial experimental expansion.

In response to your guidance, we conducted an exhaustive revision across the manuscript. In accordance with the revision guidelines, \textbf{all textual modifications, newly derived mathematical analyses, and expanded experimental discussions in the revised dissertation are highlighted in \textcolor{red}{red}} (\texttt{\textbackslash textcolor\{red\}\{...\}}).

Below, we provide a point-by-point response to your General Assessment, all 17 Major Comments, and all 11 Minor Comments and Typos.

\tableofcontents
\newpage

% ==============================================================================
\section{Part 1: Response to the General Assessment}
% ==============================================================================

\subsection{Scope, Applicative Impact, and 3D Experimental Validation}
\begin{revcomment}{General Assessment: Applicative Impact and 3D Validation}
``However, the fact that only one qualitative experiment in 3D is presented limits the thesis' potential applicative impact (...) For the audio and acoustic signal processing community, a lot remains to be explored in terms of improving and accelerating commonly used simulators.''
\end{revcomment}

\begin{response}
We agree that presenting a single qualitative 3D snapshot comparison limited the practical evaluation. In the revised dissertation, we expanded the ARD numerical experiments and the Chapter~4 limitations discussion:
\begin{enumerate}
    \item \textbf{Multi-solver 3D comparison:} The stepped enclosure has overall extents $[0,30]\times[0,30]\times[0,10]$\,m and volume $7600$\,m$^3$. The lossless frequency-response comparison includes five solver/resolution configurations: coarse and fine FDTD ($\Delta h=0.20$ and $0.10$\,m), coarse and fine TD-FEM at the same resolutions, and coarse three-subdomain ARD-PSTD.
    \item \textbf{Measured execution times:} For the reported $T=0.24$\,s runs, eight-core times are $22.1$\,s for coarse FDTD, $92.6$\,s for fine FDTD, $51.6$\,s for coarse TD-FEM, $144.3$\,s for fine TD-FEM, and approximately $36.0$\,s for ARD-PSTD. These are measurements for the stated hardware and configurations.
    \item \textbf{Separate damping comparison:} ARD-PSTD snapshots compare four values: $\nu=0$, $2.4\times10^{-6}$, $1.2\times10^{-5}$, and $5.9\times10^{-5}$\,m$^2$/s. The text describes the visible changes qualitatively over the simulated interval; the thesis does not present these as a calibrated atmospheric-absorption study.
    \item \textbf{Scope and limitations:} The cross-solver frequency-response comparison is lossless and focuses below $200$\,Hz. The 3D experiment uses a stepped axis-aligned enclosure with rigid exterior walls; it does not establish accuracy against an analytical or experimental reference, nor validate curved geometries or a broader audio-frequency range.
\end{enumerate}
\end{response}

\subsection{Analysis, Insights, and Commentary on Numerical Results}
\begin{revcomment}{General Assessment: Commentary on Damped Results}
``Unfortunately, there is very little comments or analysis of the results by the author, beyond presenting them (...) For the two damped scenarios (the most novel contributions), there is in fact no comments at all by the author, which is a problematic omission.''
\end{revcomment}

\begin{response}
We resolved this omission. In Chapter~2 (Section~2.6, Numerical Results) and Chapter~3 (Section~3.3), we added thorough physical and mathematical discussions:
\begin{enumerate}
    \item \textbf{Viscoelastic effect in the tested FDTD benchmark (Table~2.2):} For Experiment~4 (damped non-smooth pulse), the measured rate is approximately $2.22$, compared with $1.03$ in the undamped Experiment~2. This observation is consistent with viscous attenuation of high-frequency modes ($s_- \sim -\nu k_m^2$), but does not establish a general convergence-rate result.
    \item \textbf{Low-Frequency Attenuation Expansion (Section~2.1):} We derived the exact low-frequency attenuation $\alpha(\omega) \sim \frac{\nu}{2c^3}\omega^2$ for pure viscoelastic damping ($\gamma = 0$), capturing thermoviscous dissipation. For combined models ($\gamma > 0$), we derived the diffusive root $\alpha(\omega) \sim \sqrt{\frac{\gamma\omega}{2c^2}}$ and the weak-damping asymptotic $\alpha(\omega) \sim \frac{\gamma}{2c} + \frac{\nu}{2c^3}\omega^2$.
    \item \textbf{High-Frequency Modal Asymptotics (Section~2.2.4):} We corrected the modal root scaling for $k_m \to \infty$: rather than uniform $\nu k_m^2$ decay across both branches, the system becomes overdamped with two distinct branches: $s_+ \sim -c^2/\nu$ and $s_- \sim -\nu k_m^2$.
\end{enumerate}
\end{response}

\subsection{Chapter 4 Conclusion and Balance of Scientific Claims}
\begin{revcomment}{General Assessment: Concluding Chapter Balance}
``Finally, the conclusion is rather succinct and only 3 lines on future work is an order of magnitude less than what one would typically expect for a 3-year-long doctoral work. Some claims are too strong and should be qualified or retracted.''
\end{revcomment}

\begin{response}
Chapter~4 was expanded and reorganized:
\begin{enumerate}
    \item \textbf{Summary of Main Contributions (Section~4.1):} Structured as continuous academic prose (free of bulleted lists and em dashes), synthesizing FDTD vs.\ PSTD analysis, OSWR domain decomposition, ARD for viscoelastic media, and differentiable physics inversion.
    \item \textbf{Limitations of the Present Study (Section~4.2):} The discussion covers modal truncation and damped splitting accuracy, the tested geometry and PML scope, finite-grid ARD stability estimates and sampled-frequency OSWR optimization, and the limited DPNN profiles, sensors, and noise settings.
    \item \textbf{Expanded Future Research Directions (Section~4.3):} Expanded from 3 lines to several pages of detailed technical discussion across five subfields: geometric discretization (4.3.1), energy stability and exact exponential Kick splitting for ARD (4.3.2), analytical OSWR min-max theory (4.3.3), 3D scalable DPNN inversion (4.3.4), and frequency-dependent digital impedance boundary conditions and CPML (4.3.5).
\end{enumerate}
\end{response}

\newpage
% ==============================================================================
\section{Part 2: Point-by-Point Response to Major Comments}
% ==============================================================================

\begin{revcomment}{Comment R2C1: Definition of Specialized Technical Terms}
Sec. 1.3 and 1.4: Highly specialized terms are repeatedly used without defining or explaining them (e.g., ``Runge-Kutta scheme''). Even though the presented material is of a textbook nature, someone outside of the field of numerical schemes could not easily decipher it due to its density.
\end{revcomment}

\begin{response}
In Sections~1.3 and~1.4, we added explanations at first use for the numerical concepts needed to follow the discussion, including Runge--Kutta methods, the CFL condition, preconditioning, Krylov subspaces, and transmission conditions. The revision does not claim to define every specialized term in the field.
\end{response}

\begin{revcomment}{Comment R2C2: Numerical Results on Damped Benchmarks (Pages 66--67)}
For the two damped scenarios (the most novel contributions), there is in fact no comments at all by the author, which is a problematic omission.
\end{revcomment}

\begin{response}
We added in-depth discussions for Experiments~4 and~6 in Section~2.6 (Numerical Results):
\begin{itemize}
    \item \emph{Experiment~4 (Viscoelastic non-smooth pulse):} For the tested damping parameters and observation time, the measured FDTD convergence rate is approximately $2.22$, compared with $1.03$ in the undamped experiment. This numerical observation is consistent with attenuation of high-frequency content by the viscous term ($s_- \sim -\nu k_m^2$); it is not asserted as a general convergence-rate result.
    \item \emph{Experiment~6 (undamped triangular propagating pulse):} The revised text reports a PSTD error of approximately $8.19\times10^{-5}$ at $N=800$, about $42.2$ times lower than FDTD for this test. It describes the FDTD trailing ripples and the lower PSTD errors over the tested resolutions; it does not claim that PSTD error is confined strictly to boundary stencils.
\end{itemize}
\end{response}

\begin{revcomment}{Comment R2C3: Motivation for Two Overlapping Subdomains in Sec. 3.1.4}
Sec. 3.1.4: What is the motivation of using only 2 overlapping subdomains in 1D? Why not 3 or more, and why not non-overlapping?
\end{revcomment}

\begin{response}
In Section~3.1.1, we clarified that the two-subdomain overlapping geometry $\Omega_1 = (0, b)$ and $\Omega_2 = (a, L)$, with $0 < a < b < L$, is the canonical theoretical setup for deriving a frequency-dependent transmission factor \citep{gander_optimized_2007,dolean_introduction_2015,quarteroni_domain_1999}. It isolates the transmission-operator effect without introducing additional interfaces or partition-graph choices; it is not a comparison of the optimal number of subdomains. The revised text now states that three-or-more-subdomain layouts require analysis of additional interface coupling, while non-overlapping formulations require a separate analysis of their transmission conditions. We do not claim to have proved that all non-overlapping Dirichlet variants stagnate.
\end{response}

\begin{revcomment}{Comment R2C4: Retraction of Claim that FDTD is Presented for ARD}
P86: ``We present the ARD method for both FDTD and PSTD, which share the same formulation'': No, FDTD is not presented.
\end{revcomment}

\begin{response}
We clarified the scope in the Introduction, Chapter~3 overview, and Chapter~4 summary: the ARD solver implemented and evaluated here uses PSTD within the subdomains. The FDTD subsection relates undamped leapfrog to velocity Verlet as context; it does not present an implemented or evaluated FDTD-based ARD solver.
\end{response}

\begin{revcomment}{Comment R2C5: Justification of Parabolic Stability Limit}
Sec. 3.2.3: ``the second term introduces a strict parabolic stability limit ($\Delta t \le \Delta x^2/\nu$).'' This is not sufficiently detailed or justified.
\end{revcomment}

\begin{response}
We withdrew the general claim that the implemented method has the parabolic bound $\Delta t\le\Delta x^2/\nu$. The revised text reports the tested residual spectra: they were real to numerical precision but non-positive only for order 2; orders 4, 6, and 8 had small positive eigenvalues. We also added a proposition in the ARD stability analysis, proving a sufficient power-boundedness condition for the full unforced, undamped order-2 KDK update on the specified one-dimensional, cell-centered grid with two equal Neumann subdomains when $\mathrm{CFL}\le1/2$, apart from the physical mean drift. That conservative result does not cover the damped or higher-order updates. The reported full-matrix thresholds remain numerical estimates and do not prove a general stability limit.
\end{response}

\begin{revcomment}{Comment R2C6: Initialization of the Iterative ARD Algorithm}
Sec. 3.2.4: How is the presented iterative algorithm initialised, i.e., what are the initial values of $U_m^n$ and $V_m^n$?
\end{revcomment}

\begin{response}
In Section~3.2.4, we stated that the initial modal states are obtained by applying the forward DCT to the sampled initial displacement and velocity: $\mathbf U_m^0=\mathcal D\{\mathbf u_0\}$ and $\mathbf V_m^0=\mathcal D\{\mathbf v_0\}$. 
\end{response}

\begin{revcomment}{Comment R2C7: Softening ``Validating 3D ARD'' and Adding Quantitative Evaluation}
P96: ``validating the accuracy of the multi-domain pseudo-spectral time-domain (PSTD) formulation in 3D.'' This is too strong a claim. A single configuration is tested, with two (very close) values of the viscoelastic parameter. Only qualitative snapshots are displayed.
\end{revcomment}

\begin{response}
The revised thesis presents a numerical study rather than a validation of general 3D accuracy. It compares five solver/resolution configurations in a lossless frequency-response analysis, reports runtimes for the stated setup, and separately shows ARD-PSTD snapshots for four damping values. The damping sweep is not a five-solver comparison, and neither experiment provides an analytical or experimental ground truth.
\end{response}

\begin{revcomment}{Comment R2C8: Internal Boundary Handling in 3D ARD Snapshots}
Fig. 3.12: ``The solver seamlessly handles the internal boundaries between the three rectangular subdomains.'' One cannot actually see that from the figure.
\end{revcomment}

\begin{response}
Figure~3.11 was replaced with high-resolution snapshots showing the wavefield across the virtual interface lines $\Gamma_1$ and $\Gamma_2$, marked with dashed red lines. The snapshots indicate visually continuous fields in the displayed case; they do not by themselves establish the absence of reflections or phase distortions.
\end{response}

\begin{revcomment}{Comment R2C9: 3D Damping Visualization and Physical Scaling}
Fig. 3.16 is not very informative. It is impossible to see differences between the 3 images. Plotting the error or relative error would be better.
\end{revcomment}

\begin{response}
We agree that the earlier plots did not make quantitative differences easy to assess. In the revised DPNN figures, we added field-error panels for the displayed cases. For the 3D ARD experiment, we retained wavefield snapshots and added a shared colorbar and four damping values; these remain qualitative comparisons rather than relative-error plots.
\begin{enumerate}
    \item \textbf{DPNN error panels in Figures~3.18, 3.21, and 3.24:} The revised spatiotemporal figures show the ground-truth field alongside absolute-error panels for FDTD-DPNN, PSTD-DPNN, and the PINN baseline, using a shared scale.
    \begin{itemize}
        \item Panel~(a): Ground-truth wavefield $u^*(x,t)$ (seismic bipolar colormap).
        \item Panel~(b): FDTD-DPNN absolute error $|u_{\mathrm{FDTD}}(x,t) - u^*(x,t)|$.
        \item Panel~(c): PSTD-DPNN absolute error $|u_{\mathrm{PSTD}}(x,t) - u^*(x,t)|$.
        \item Panel~(d): Standard PINN baseline absolute error $|u_{\mathrm{PINN}}(x,t) - u^*(x,t)|$.
    \end{itemize}
    The panels show the spatial error patterns in the displayed cases. Quantitative comparisons in Table~3.2 show that performance depends on the initial profile, with the PINN baseline achieving lower field error on the multimodal standing-wave case.
    \item \textbf{3D damping comparisons:} We added ARD-PSTD snapshots for four values of $\nu$ and described the visible changes over the simulated interval. The thesis treats this as a qualitative comparison; it does not claim a frequency-domain measurement of attenuation versus $\nu$.
\end{enumerate}
\end{response}

\begin{revcomment}{Comment R2C10: Explicit Description of Cosine Annealing in DPNN}
The loss plots (Fig. 3.14, 3.17, 3.20...) feature abrupt changes, suggesting that some elaborate techniques are used to modify the learning rate along the way, but there is no explicit mention of this despite a mention of ``cosine annealing'' which is not described at all.
\end{revcomment}

\begin{response}
After comparing the text with the checked-in benchmark script, we removed the unsupported cosine-annealing description. The script uses two Adam optimizers with fixed learning rates, $10^{-3}$ for network weights and $5\times10^{-2}$ for the logarithmic physical parameters; it has no scheduler. It trains for 350 epochs.
\begin{itemize}
    \item Network weights $\Theta$ use the fixed learning rate $\eta=10^{-3}$.
    \item The logarithmic parameters $(\xi_\gamma,\xi_\nu)$ use the fixed learning rate $\eta=5\cdot10^{-2}$.
\end{itemize}
\end{response}

\begin{revcomment}{Comment R2C11: Critical Scope and Limitations of DPNN Experiments}
Sec 3.3.1: As in chapter 3, the comments on the experiments are extremely limited. This section also severely lacks quantitative analysis. There are 3 qualitative experiments, which seems too limited to rigorously validate the approach. Especially since quite strong limitations seem to appear in the third experiment.
\end{revcomment}

\begin{response}
We thank the Reviewer for this observation. In the revised thesis, we expanded Section~3.3 with quantitative benchmarks, a high-capacity baseline comparison, disentangled inversions, and regularization ablations. These experiments broaden the evidence for the selected one-dimensional configurations, while the limitations of this scope remain explicitly acknowledged:
\begin{enumerate}
    \item \textbf{Quantitative benchmark comparison and PINN baseline (Section~3.3.2, Tables~3.2 and~3.3):} We added a numerical comparison reporting five quantities for the three tested initial-condition benchmarks (multimodal standing waves, smooth Gaussian pulses, and non-smooth triangular pulses), comparing PSTD-DPNN and FDTD-DPNN with one PINN baseline \citep{raissi_pinn_2019,haghighat_sciann_2021}:
    \begin{itemize}
        \item Normalized Mean Squared Error on initial pressure: $\mathrm{NMSE}(p_0) = \|p_0^{\mathrm{pred}} - p_0^*\|_2^2 / \|p_0^*\|_2^2$.
        \item Full spatiotemporal wavefield relative $L_2$ error: $\|u_{\mathrm{pred}} - u^*\|_2 / \|u^*\|_2$ (\%).
        \item Physical attenuation parameter estimation: telegrapher damping $\hat{\gamma}$ and viscoelastic damping $\hat{\nu}$ alongside relative errors.
        \item Wall-clock training time per epoch (ms/epoch). The checked-in script explicitly runs on CPU; it does not use the available GPU, and does not record the host CPU model.
        \item \emph{Results in these configurations}: On the tested triangular pulse, PSTD-DPNN attains $\mathrm{NMSE}=4.49\times10^{-3}$ and field relative error $3.09\%$, while FDTD-DPNN and the PINN baseline yield $14.67\%$ and $28.20\%$ field error. The PINN baseline also attains $2.80\%$ field error on the tested standing-wave case. PSTD-DPNN requires 12.1--13.8\,ms/epoch compared with 91.3--95.1\,ms/epoch for this PINN setup, which uses 300 interior collocation points per epoch. These measurements compare the selected implementations and settings; they do not establish superiority over PINNs or black-box inverse methods in general.
    \end{itemize}
    \item \textbf{Disentangled Inversion Experiments (Section~3.3.3, Table~3.4):} To compare parameter-only and joint inversion in one Gaussian benchmark, we report three settings:
    \begin{itemize}
        \item \emph{Experiment A ($u_0$ reconstruction only, $\gamma, \nu$ known)}: Attains an initial condition NMSE of $3.23 \times 10^{-3}$ and a field error of $3.94\%$.
        \item \emph{Experiment B (parameter identification only, $u_0$ known)}: For this Gaussian benchmark and initialization, logarithmic parameterization identifies $\nu$ with $3.0\%$ relative error ($\hat{\nu} = 4.85 \times 10^{-3}\,\mathrm{m}^2/\mathrm{s}$) and $\gamma$ with $10.8\%$ relative error ($\hat{\gamma} = 0.1108\,\mathrm{s}^{-1}$), with field error $0.37\%$. This is a result for the tested configuration and does not establish general parameter-recovery accuracy.
        \item \emph{Experiment C (Joint inversion, $u_0, \gamma, \nu$ unknown)}: Directly aligns with the main Gaussian benchmark (Table~3.2), achieving $1.68\%$ field error (NMSE $3.00 \times 10^{-3}$), recovering $\nu$ within $2.0\%$ ($\hat{\nu} = 5.10 \times 10^{-3}\,\mathrm{m}^2/\mathrm{s}$) and $\gamma$ within $3.5\%$ ($\hat{\gamma} = 0.0965\,\mathrm{s}^{-1}$).
    \end{itemize}
    \item \textbf{Regularization Ablation Study and Elimination of $H^1$ (Section~3.3.4 \& Table~3.5):} We conducted a systematic 5-configuration ablation study (Data only, $L_1$ only, $H^1$ only, Coupled $L_1+H^1$, Decoupled $L_1+H^1$):
    \begin{itemize}
        \item \emph{Data fitting without explicit spatial regularization}: Among the evaluated configurations, data-only fitting gives the better combined field-and-parameter trade-off; $L_1$ lowers Gaussian initial-condition NMSE but worsens its field error.
        \item \emph{Tested Sobolev configurations}: The $H^1$ penalty suppresses high-wavenumber content, which can remove information relevant to estimating physical dissipation. In the tested coupled configuration the estimated $\nu$ is very small; the ablation does not establish this mechanism for other settings.
    \end{itemize}
\end{enumerate}
The findings are synthesized in the Concluding Assessment (Section~3.3.5).
\end{response}

\begin{revcomment}{Comment R2C12: Retraction of ``Remarkable Robustness'' on Triangular Pulse}
P104: ``Figures 3.20 to 3.22 showcase the remarkable robustness of the carefully calibrated Sobolev H1 regularization''. To be honest, not really. The relative reconstruction error on the initial condition seems very large. This claim should be retracted or further expanded.
\end{revcomment}

\begin{response}
We retracted the phrase ``remarkable robustness''. Following your critique, our new systematic ablation study (Section~3.3.4, Table~3.5) revealed that Sobolev $H^1$ regularization was actually detrimental rather than robust for this non-smooth profile:
\begin{itemize}
    \item The proposed unregularized formulation achieves an initial condition NMSE of only $4.49 \times 10^{-3}$ and a full-field relative error of $3.09\%$, recovering $\gamma$ within $4.7\%$ and $\nu$ within $3.6\%$.
    \item In contrast, imposing Sobolev $H^1$ regularization degrades field error to $18.10\%$ and causes estimated physical viscosity $\nu$ to collapse to $7.66 \times 10^{-6}$ ($99.85\%$ relative error).
    \item The thesis notes that the $H^1$ penalty suppresses high-wavenumber content and can remove information relevant to estimating physical dissipation. In the tested coupled-weight configuration, the estimated viscosity is very small; the ablation does not establish this mechanism for other settings.
\end{itemize}
\end{response}

\begin{revcomment}{Comment R2C13: Toning Down ``Unprecedented Computational Efficiency''}
P107: ``achieve unprecedented computational efficiency and spatial accuracy.'' There is no solid proof of this in the thesis, this claim should be tuned down or retracted.
\end{revcomment}

\begin{response}
We retracted the superlative claim. In the reported 3D setup, eight-core ARD-PSTD runtime is about $36.0$\,s, compared with $22.1$\,s for coarse FDTD, $92.6$\,s for fine FDTD, and $144.3$\,s for fine TD-FEM. The lossless frequency-response comparison shows closer agreement below $200$\,Hz for the tested fine-grid references; it does not establish zero dispersion or absolute accuracy.
\end{response}

\begin{revcomment}{Comment R2C14: Retraction of Claim Extending ARD to FDTD}
P107: ``Subsequently, we extended the Adaptive Rectangular Decomposition (ARD) method to both FDTD and PSTD frameworks.'': No, this is incorrect, the FDTD framework is not treated. The author mentions it in a paragraph, but then explicitly writes that it won't be treated.
\end{revcomment}

\begin{response}
We clarified the scope in the Introduction, Chapter~3 overview, and Chapter~4 summary: the ARD solver implemented and evaluated here uses PSTD within the subdomains. The FDTD subsection relates undamped leapfrog to velocity Verlet as context; it does not present an implemented or evaluated FDTD-based ARD solver.
\end{response}

\begin{revcomment}{Comment R2C15: Toning Down Superlative Claims on Sobolev Regularization}
P108: ``(...) Sobolev $H_1$ smoothness penalty gracefully suppressed high-frequency artifacts without smearing out sharp discontinuities'': This is a highly subjective and superlative statement. In fact, there is no real rigorous treatment of this, just a succinct qualitative experiment.
\end{revcomment}

\begin{response}
We retracted ``gracefully suppressed (...) without smearing out sharp discontinuities''. We replaced it with a five-configuration numerical ablation in Section~3.3.4 (Table~3.5) and the concluding assessment in Section~3.3.5:
\begin{itemize}
    \item In the tested triangular case, the $H^1$ configurations have larger reconstruction errors and do not preserve the sharp profile as well as data-only fitting.
    \item The reported field errors for $H^1$-only fitting are $7.42\%$ on the Gaussian profile and $18.10\%$ on the triangular profile; the tested coupled configuration also returns a very small viscosity estimate. These observations are limited to the profiles and penalty weights in the ablation.
    \item For the Gaussian profile, $L^1$ lowers the initial-condition NMSE but increases field error to $2.42\%$ from $1.68\%$. For the triangular profile, data-only fitting gives the lowest reported field error ($3.09\%$) among the tested configurations. The thesis therefore does not claim that explicit regularization is generally beneficial or generally harmful.
\end{itemize}
\end{response}

\begin{revcomment}{Comment R2C16: Retraction of Unbacked Claim on Black-Box Inversions}
P108: ``(...), significantly outperforming purely data-driven black-box inversions'': this claim has no supporting evidence and should be retracted. No comparison is provided for this method, nor any quantitative results.
\end{revcomment}

\begin{response}
We retracted the unbacked phrase ``significantly outperforming purely data-driven black-box inversions''. Table~3.2 reports quantitative results for the tested differentiable-physics solvers and one PINN baseline (4 hidden layers, 128 units, $50{,}049$ network parameters plus two damping variables) \citep{raissi_pinn_2019,haghighat_sciann_2021}. In the tested non-smooth case, PSTD-DPNN has lower reconstruction error than this PINN setup, and the measured epoch times are lower in the CPU-only script configuration:
\begin{itemize}
    \item On the tested triangular pulse, the reported field errors are $3.09\%$ for PSTD-DPNN and $28.20\%$ for this PINN baseline.
    \item Under the reported implementation and CPU-only configuration, PSTD-DPNN's measured epoch time is 6.6 to 7.7 times lower than this PINN setup. This comparison is configuration-dependent and does not imply the same speed ratio for other architectures or implementations.
    \item The DPNN SIREN has 2,209 weights and biases; including its two damping variables gives 2,211 trainable quantities. The PINN has 50,049 network weights and biases plus two damping variables. This counts the specific networks in the checked-in benchmark script.
\end{itemize}
\end{response}

\begin{revcomment}{Comment R2C17: Expansion of Future Research Directions}
P108: 3 lines on future work is an order of magnitude less than what one would typically expect for a 3-year-long doctoral work.
\end{revcomment}

\begin{response}
Section~4.3 was expanded from 3 lines to over 4 pages of continuous technical prose, detailing 5 distinct research avenues:
\begin{enumerate}
    \item \emph{Section~4.3.1:} Geometric discretization and staircasing boundary error mitigation (Cartesian cut-cell, embedded boundary, curvilinear transformations).
    \item \emph{Section~4.3.2:} Energy stability and higher-order time integration for ARD (exact exponential Kick flow, midpoint Runge-Kutta, discrete energy estimates, high-order DtN Padé operators).
    \item \emph{Section~4.3.3:} Analytical min-max theory and adaptive transmission for OSWR.
    \item \emph{Section~4.3.4:} Scaling DPNN to 3D and advanced inversion formulations (3D scaling, optimal sensor placement, edge-preserving Total Variation regularization, and comparison with FWI and TRA).
    \item \emph{Section~4.3.5:} Realistic acoustic modelling, advanced boundary conditions (CPML/CFS-PML with auxiliary memory variables, digital impedance filters), and GPU acceleration.
\end{enumerate}
\end{response}

\newpage
% ==============================================================================
\section{Part 3: Point-by-Point Response to Minor Comments and Typos}
% ==============================================================================

\begin{revcomment}{Comment R2m1: Timeline of Wave-Based Methods (P4)}
The author writes that ``The first time-dependent wave-based methods appeared around the same period of time [(1998)]'', but after that, papers from 1968, 1971 and 1974 are cited. Did time-dependent wave-based methods appear in the late 90s or early 70s? Please clarify.
\end{revcomment}

\begin{response}
We clarified Section~1.1 to distinguish the fields and timelines: the current text says that time-dependent wave-based methods date to the late 1960s and 1970s in geophysics and seismology \citep{alterman_propagation_1968,ottaviani_elastic-wave_1971,alford_accuracy_1974}, while their application to room acoustics emerged in the 1990s \citep{botteldooren_finitedifference_1995}. It no longer attributes the earlier developments to room acoustics.
\end{response}

\begin{revcomment}{Comment R2m2: Structural Placement of DPNN Paragraphs (P12)}
The three paragraphs on DPNNs talk about inverse problem while the 5 previous pages in the section are on numerical schemes for forward problems, so they seem a bit out of place. This should maybe be a different section.
\end{revcomment}

\begin{response}
We restructured Section~1.4 into dedicated, clearly demarcated subsections: Section~1.4.1 details forward numerical solvers, while Section~1.4.2 formally introduces inverse acoustic problems and motivates Differentiable Physics Neural Networks.
\end{response}

\begin{revcomment}{Comment R2m3: Equation Reference for $k(i\omega)$ (P18)}
``$k(i\omega)$ satisfies the dispersion relation (2.2).'': I think this is the wrong equation reference. This should probably point to the equation in section 2.1 (which doesn't have a number).
\end{revcomment}

\begin{response}
We assigned an equation label to the continuous dispersion relation in Section~2.1 and corrected the cross-reference to point directly to Equation~(2.1).
\end{response}

\begin{revcomment}{Comment R2m4: Clarification of the Three Models (P24)}
``An important qualitative difference between the three models considered...'' I am confused what are the three models referred to here. Section 2.1 only mentions two: multi-relaxation and telegrapher + viscoelastic damping. Is the undamped model considered as a third model?
\end{revcomment}

\begin{response}
We clarified the text: the three physical models compared in Section~2.1 are (1) the lossless undamped wave equation ($\gamma = \nu = 0$), (2) the telegrapher model ($\gamma > 0, \nu = 0$), and (3) the combined telegrapher and viscoelastic model ($\gamma > 0, \nu > 0$).
\end{response}

\begin{revcomment}{Comment R2m5: Role of Proposition 2.4 (P40)}
The usefulness of stating Proposition 2.4 in full detail is not clear, as this seems completely standard.
\end{revcomment}

\begin{response}
We clarified the scope of Proposition~2.4 in the text immediately following it. The identity $T_m=A_m^{-1}(S_m-I)$ is included as an implementation formula for the piecewise-constant modal forcing update on non-constant modes; the singular constant mode is evaluated directly in Appendix~A. We do not present this standard identity as a new result or as a stability theorem.
\end{response}

\begin{revcomment}{Comment R2m6: Clarification of ``Advanced in the Modal Domain'' (P54)}
``The numerical solution is advanced in the modal domain.'' I don't understand this sentence.
\end{revcomment}

\begin{response}
We replaced this sentence with a clear operational description: \emph{``The spatial acoustic field is transformed into discrete modal coefficients via the DCT-II, advanced in time analytically via decoupled modal transition operators, and synthesized back into physical space via the inverse DCT-II.''}
\end{response}

\begin{revcomment}{Comment R2m7: Grammatical Structure in Section 3.1 (P71)}
``For the classical Schwarz Waveform Relaxation method, for which Dirichlet transmission conditions are employed, whereas for Optimized Schwarz Waveform Relaxation...'' There appears to be a structure problem with this sentence (remove ``for which''?).
\end{revcomment}

\begin{response}
We corrected the sentence syntax by removing the redundant relative clause: \emph{``In the classical Schwarz Waveform Relaxation method, Dirichlet transmission conditions are employed, whereas Optimized Schwarz Waveform Relaxation utilizes generalized Robin transmission conditions.''}
\end{response}

\begin{revcomment}{Comment R2m8: Overlap Definition $a < b$ (P81)}
In the definition of $\Omega_1$ and $\Omega_2$: I assume that $a < b$ since the domains are overlapping? Please specify.
\end{revcomment}

\begin{response}
In Section~3.1.1, we explicitly stated that $0 < a < b < L$, defining the positive overlap length as $\delta = b - a > 0$.
\end{response}

\begin{revcomment}{Comment R2m9: Duplicate Citation Typo (P86)}
``(Strang Strang (1968))''
\end{revcomment}

\begin{response}
We fixed the citation command in Section~3.2 to use \texttt{\textbackslash citep\{strang\_construction\_1968\}}, eliminating the duplicate author name.
\end{response}

\begin{revcomment}{Comment R2m10: Citation Typo (P89)}
``Strang splitting Strang (1968)'': typo (one should use \texttt{\textbackslash citep} here).
\end{revcomment}

\begin{response}
We corrected the inline citation to standard parenthetical format: \emph{``Strang splitting \citep{strang_construction_1968}''}.
\end{response}

\begin{revcomment}{Comment R2m11: Definition of Drift and Kick Operators (P90)}
``$D(\Delta t)$ denotes the exact modal propagator and $K(\Delta t/2)$ represents...'' $D$ and $K$ are not adequately defined here.
\end{revcomment}

\begin{response}
In Section~3.2.2, we provided formal mathematical definitions: $D(\Delta t) = \mathcal{S}_{\text{local}}(\Delta t)$ denotes the exact modal time-evolution operator for the uncoupled subdomain interior, while $K(\tau) = I + \tau B$ denotes the discrete Kick operator integrating the singular interface transmission operator.
\end{response}

\vspace{0.8cm}
\noindent\rule{\textwidth}{0.4pt}\\[0.5em]
We hope these revisions and the added numerical results address the concerns raised. We thank you once again for your constructive guidance.

\vspace{0.5cm}
\noindent Sincerely,\\[0.3em]
\textbf{Gerardo Cicalese}\\
\emph{Ph.D. Candidate in Information Technology, Politecnico di Milano}\\[0.5em]
\textbf{Prof.~Ilario Mazzieri}\\
\emph{Advisor, Associate Professor, Politecnico di Milano}\\[0.5em]
\textbf{Prof.~Gabriele Ciaramella}\\
\emph{Co-Advisor, Associate Professor, Politecnico di Milano}\\[0.5em]
\textbf{Prof.~Fabio Antonacci}\\
\emph{Co-Advisor, Associate Professor, Politecnico di Milano}

\newpage
\bibliographystyle{unsrtnat-nourl}
\bibliography{Thesis_bibliography}

\end{document}
